\clearpage\corrigesChapitre{12}{Forces et interactions}
\corrigesSouples
\begin{blocLivre}{cor-12-1}\begin{corrige}{1}{Nommer l’acteur}
\begin{questions}
\item La valise reçoit l’action de la main.
\item $\vect F_{\mathrm{main}\to\mathrm{valise}}$, appliquée au contact de la poignée.
\end{questions}
\end{corrige}\end{blocLivre}
\begin{blocLivre}{cor-12-2}\begin{corrige}{2}{Poids d’un objet}
\begin{questions}
\item $P=0,200\times10,0=\qty{2,00}{N}$.
\item Verticale locale, vers le bas, c’est-à-dire vers le centre de la Terre.
\end{questions}
\end{corrige}\end{blocLivre}
\begin{blocLivre}{cor-12-3}\begin{corrige}{3}{Contact ou distance}
\schema[corrige]{bilan-forces}
\begin{questions}
\item Action gravitationnelle de la Terre et action de la table.
\item Terre : à distance ; table : contact. On néglige ici l’action de l’air.
\end{questions}
\end{corrige}\end{blocLivre}
\begin{blocLivre}{cor-12-4}\begin{corrige}{4}{Deux patineurs}
\schema[corrige]{reciprocite}
\begin{questions}
\item Les valeurs sont égales et les sens opposés, selon les actions réciproques.
\item Chaque force agit sur un patineur différent.
\end{questions}
\end{corrige}\end{blocLivre}
\begin{blocLivre}{cor-12-5}\begin{corrige}{5}{Changer de planète}
\begin{questions}
\item $P_T=29,4$ N et $P_L=4,8$ N, avant arrondi adapté.
\item La masse de 3,0 kg reste la même ; le champ de pesanteur local et le poids changent.
\end{questions}
\end{corrige}\end{blocLivre}
\begin{blocLivre}{cor-12-6}\begin{corrige}{6}{Une attraction gravitationnelle}
\begin{questions}
\item $F=6,67\times10^{-11}\times10,0\times20,0/0,50^2=\qty{5,3e-8}{N}$ à deux chiffres significatifs.
\item Elle est divisée par quatre, car $F$ dépend de $1/d^2$.
\end{questions}
\end{corrige}\end{blocLivre}
\begin{blocLivre}{cor-12-7}\begin{corrige}{7}{Une boule suspendue}
\schema[corrige]{suspension}
\begin{questions}
\item Poids vertical vers le bas ; traction du fil vers le haut. La figure montre les deux vecteurs.
\item $P=2,00$ N ; l’équilibre donne une tension de 2,00 N.
\end{questions}
\end{corrige}\end{blocLivre}
\begin{blocLivre}{cor-12-8}\begin{corrige}{8}{Le centre compte}
\begin{questions}
\item $d=6400+400=6800$ km, soit $6,8\times10^6$ m.
\item La distance de la loi est celle entre les centres, non l’altitude au-dessus de la surface.
\end{questions}
\end{corrige}\end{blocLivre}
\clearpage
\begin{blocLivre}{cor-12-9}\begin{corrige}{9}{Terre et pomme}
\begin{questions}
\item Une attraction de 1,5 N, de sens opposé, dirigée de la Terre vers la pomme.
\item Non. L’égalité des valeurs ne dépend pas de l’égalité des masses ; les effets sur leurs mouvements peuvent être très différents.
\end{questions}
\end{corrige}\end{blocLivre}
\begin{blocLivre}{cor-12-10}\begin{corrige}{10}{Un dessin à l’échelle}
\begin{questions}
\item $4,5/1,5=\qty{3,0}{cm}$.
\item Le point d’application, la direction et le sens, en plus de la valeur et de l’échelle.
\end{questions}
\end{corrige}\end{blocLivre}
\begin{blocLivre}{cor-12-11}\begin{corrige}{11}{Un support incliné}
\begin{questions}
\item Le support exerce une composante normale et une composante tangentielle de frottement.
\item Vers le haut du plan, pour s’opposer au glissement envisagé.
\end{questions}
\end{corrige}\end{blocLivre}
\begin{blocLivre}{cor-12-12}\begin{corrige}{12}{Retrouver g}
\demarche{Identifier le poids à l’attraction gravitationnelle à la surface de la planète.}
\begin{questions}
\item $g=GM/R^2$ après simplification par $m$.
\item $g\simeq\qty{9,8}{N.kg^{-1}}$. La masse de l’objet n’apparaît plus dans cette expression.
\end{questions}
\retiens{La masse de l’objet se simplifie dans le calcul de g ; celle de la planète reste présente.}
\end{corrige}\end{blocLivre}

